\section{Cote 18, physical 49/footer 48: simultaneous local genericity}

\subsection{The original assertion and its two hypotheses}
The bottom of the original sheet, cote 18, physical PDF page 49,
manuscript/footer 48, considers a connected affine algebraic group
\(G/\mathbf Q\), a compact open subgroup
\(U\subset G(\mathbf A_f)\), and a coset \(P=Ua\).
It asserts the existence in \(P\) of an element whose component at
every finite prime is generic over \(\mathbf Q\).  The preceding
lines apply this to tuples of comparison matrices; physical page 51,
footer 50, continues the assertion in terms of the local points of
the variety of tuples.  Here \(\mathbf A_f\) denotes the finite
adeles, written with the archive's hatted rational notation there.

The group is explicitly connected on this sheet.  The proposed
motivic application separately assumes a uniform form of the Tate
conjecture giving a finite-index adelic Galois image.  The local
genericity assertion below is unconditional.  Its application to a
specified Galois representation retains the compact-open-image
hypothesis; the adopted RH and Dirichlet GRH are not used to supply it.

\subsection{BR-18-AD01: the adelic assertion, including finite tuples}
Let \(G\) be a connected affine algebraic group over \(\mathbf Q\),
let \(r\geq1\), and let
\(\mathcal O\subset G(\mathbf A_f)^r\) be a nonempty open set.
There exists \(g=(g_1,\ldots,g_r)\in\mathcal O\) such that,
for every rational prime \(\ell\), evaluation is injective:
\[
 \mathbf Q[G^r]\longrightarrow\mathbf Q_\ell,
 \qquad f\longmapsto f(g_{1,\ell},\ldots,g_{r,\ell}).
 \tag{18.AD1}
\]
The set of such tuples is dense in \(G(\mathbf A_f)^r\).
In particular it meets every coset of a compact open subgroup.
Genericity in (18.AD1) is relative to \(\mathbf Q\), exactly as in
the source; it does not mean that a \(\mathbf Q_\ell\)-valued point
is generic after allowing arbitrary \(\mathbf Q_\ell\)-coefficients.

\paragraph{Known local input.}
Poonen (2017), Proposition 3.5.75 and Remark 3.5.76, p.~87,
give the following stronger local form of Zariski density: in an
irreducible variety over a local field, every open neighborhood of
a smooth rational point contains a point outside any prescribed
countable union of proper closed subvarieties.  This is the known
local result; we give only the additional adelic deduction.
For the group hypotheses, the same source, Proposition 5.2.4,
pp.~121--122, makes the identity component geometrically
irreducible, and Corollary 5.2.18, p.~125, gives smoothness in
characteristic zero.

\paragraph{Complete proof of the deduction.}
Put \(X=G^r\).  The cited group facts imply that \(X\) and every
\(X_{\mathbf Q_\ell}\) are smooth and geometrically integral.
The coordinate ring \(A=\mathbf Q[X]\) is countable: it is a
finitely generated algebra over the countable field \(\mathbf Q\).
For each nonzero \(f\in A\), its image in
\(A\otimes_{\mathbf Q}\mathbf Q_\ell\) remains nonzero, by
faithful flatness of field extension.  Its zero set
\(Z_{f,\ell}\) is consequently a proper closed subvariety of
\(X_{\mathbf Q_\ell}\).  The local result applies to this
countable family.  Thus every nonempty analytic open subset of
\(X(\mathbf Q_\ell)\) contains a point \(x_\ell\) with
\(f(x_\ell)\ne0\) for all \(0\ne f\in A\).

Choose a point of \(\mathcal O\) and then a basic restricted-product
neighborhood of it contained in \(\mathcal O\).  Explicitly, after
choosing an affine group model of \(G\) outside finitely many
primes, this neighborhood has the form
\[
       \prod_{\ell\in S}V_\ell
          \times\prod_{\ell\notin S}\mathcal G(\mathbf Z_\ell)^r,
 \tag{18.AD2}
\]
where \(S\) is finite and every \(V_\ell\) is a nonempty
analytic open subset of \(X(\mathbf Q_\ell)\).
Each \(\mathcal G(\mathbf Z_\ell)^r\) appearing here is also
nonempty and analytically open: it contains the identity, and in
an affine embedding its integrality conditions are open conditions
in \(\mathbf Q_\ell\).  These are precisely the basic sets in
the restricted-product topology.

For each \(\ell\), use the preceding local conclusion to choose
\(x_\ell\) in the corresponding factor of (18.AD2).  The tuple
\((x_\ell)_\ell\) is an adelic point because it is integral
outside \(S\); it lies in \(\mathcal O\).  Every evaluation map
in (18.AD1) has zero kernel by construction.  Since
\(\mathcal O\) was arbitrary, the set of such points is dense.
Applying this with \(r=1\) and \(\mathcal O=Ua\) proves the
original assertion. \(\square\)

\paragraph{The precise Galois-image application.}
Suppose a group homomorphism
\(\rho:\Pi\to G(\mathbf A_f)\) has compact open image \(U\).
For prescribed \(a_1,\ldots,a_r\in G(\mathbf A_f)\), there
exist \(\sigma_1,\ldots,\sigma_r\in\Pi\) such that the tuple
\[
       (\rho(\sigma_1)a_1,\ldots,\rho(\sigma_r)a_r)
 \tag{18.AD3}
\]
satisfies (18.AD1) at every \(\ell\).  Indeed,
\(Ua_1\times\cdots\times Ua_r\) is nonempty and open; apply
the theorem and lift each of its \(U\)-coordinates through
the surjection \(\rho:\Pi\twoheadrightarrow U\).
This proves the finite-tuple step of the archive argument under
its stated image hypothesis.

It matters that this is an existence statement for arbitrary
elements of a compact open image.  Even a dense distinguished
subset need not contain a generic element.  For example,
\(\mathbf Z\subset\widehat{\mathbf Z}\subset\mathbf A_f\)
is dense in the compact open additive subgroup
\(\widehat{\mathbf Z}\): any finite list of residue conditions
has an integer solution by the Chinese remainder theorem.
Nevertheless each integer \(n\) is killed under evaluation by
the nonzero polynomial \(t-n\in\mathbf Q[t]\) at every prime.
Thus density of a selected collection of arithmetic elements
does not by itself provide the element of (18.AD1).

\paragraph{Bibliographic locator.}
Poonen, B. (2017). \emph{Rational points on varieties}.
Graduate Studies in Mathematics, 186. American Mathematical Society.
Proposition 3.5.75 and Remark 3.5.76, p.~87;
Proposition 5.2.4, pp.~121--122; Corollary 5.2.18, p.~125.
Author-hosted text:
\url{https://math.mit.edu/~poonen/papers/Qpoints.pdf}.

